\IfFileExists{stacks-project.cls}{%
\documentclass{stacks-project}
}{%
\documentclass{amsart}
}

% For dealing with references we use the comment environment
\usepackage{verbatim}
\newenvironment{reference}{\comment}{\endcomment}
%\newenvironment{reference}{}{}
\newenvironment{slogan}{\comment}{\endcomment}
\newenvironment{history}{\comment}{\endcomment}

% For commutative diagrams we use Xy-pic
\usepackage[all]{xy}

% We use 2cell for 2-commutative diagrams.
\xyoption{2cell}
\UseAllTwocells

% We use multicol for the list of chapters between chapters
\usepackage{multicol}

% This is generally recommended for better output
\usepackage{lmodern}
\usepackage[T1]{fontenc}

% For cross-file-references
\usepackage{xr-hyper}

% Package for hypertext links:
\usepackage{hyperref}

% For any local file, say "hello.tex" you want to link to please
% use \externaldocument[hello-]{hello}
\externaldocument[introduction-]{introduction}
\externaldocument[conventions-]{conventions}
\externaldocument[sets-]{sets}
\externaldocument[categories-]{categories}
\externaldocument[topology-]{topology}
\externaldocument[sheaves-]{sheaves}
\externaldocument[sites-]{sites}
\externaldocument[stacks-]{stacks}
\externaldocument[fields-]{fields}
\externaldocument[algebra-]{algebra}
\externaldocument[brauer-]{brauer}
\externaldocument[homology-]{homology}
\externaldocument[derived-]{derived}
\externaldocument[simplicial-]{simplicial}
\externaldocument[more-algebra-]{more-algebra}
\externaldocument[smoothing-]{smoothing}
\externaldocument[modules-]{modules}
\externaldocument[sites-modules-]{sites-modules}
\externaldocument[injectives-]{injectives}
\externaldocument[cohomology-]{cohomology}
\externaldocument[sites-cohomology-]{sites-cohomology}
\externaldocument[dga-]{dga}
\externaldocument[dpa-]{dpa}
\externaldocument[sdga-]{sdga}
\externaldocument[hypercovering-]{hypercovering}
\externaldocument[schemes-]{schemes}
\externaldocument[constructions-]{constructions}
\externaldocument[properties-]{properties}
\externaldocument[morphisms-]{morphisms}
\externaldocument[coherent-]{coherent}
\externaldocument[divisors-]{divisors}
\externaldocument[limits-]{limits}
\externaldocument[varieties-]{varieties}
\externaldocument[topologies-]{topologies}
\externaldocument[descent-]{descent}
\externaldocument[perfect-]{perfect}
\externaldocument[more-morphisms-]{more-morphisms}
\externaldocument[flat-]{flat}
\externaldocument[groupoids-]{groupoids}
\externaldocument[more-groupoids-]{more-groupoids}
\externaldocument[etale-]{etale}
\externaldocument[chow-]{chow}
\externaldocument[intersection-]{intersection}
\externaldocument[pic-]{pic}
\externaldocument[weil-]{weil}
\externaldocument[adequate-]{adequate}
\externaldocument[dualizing-]{dualizing}
\externaldocument[duality-]{duality}
\externaldocument[discriminant-]{discriminant}
\externaldocument[derham-]{derham}
\externaldocument[local-cohomology-]{local-cohomology}
\externaldocument[algebraization-]{algebraization}
\externaldocument[curves-]{curves}
\externaldocument[resolve-]{resolve}
\externaldocument[models-]{models}
\externaldocument[functors-]{functors}
\externaldocument[equiv-]{equiv}
\externaldocument[pione-]{pione}
\externaldocument[etale-cohomology-]{etale-cohomology}
\externaldocument[proetale-]{proetale}
\externaldocument[relative-cycles-]{relative-cycles}
\externaldocument[more-etale-]{more-etale}
\externaldocument[trace-]{trace}
\externaldocument[crystalline-]{crystalline}
\externaldocument[spaces-]{spaces}
\externaldocument[spaces-properties-]{spaces-properties}
\externaldocument[spaces-morphisms-]{spaces-morphisms}
\externaldocument[decent-spaces-]{decent-spaces}
\externaldocument[spaces-cohomology-]{spaces-cohomology}
\externaldocument[spaces-limits-]{spaces-limits}
\externaldocument[spaces-divisors-]{spaces-divisors}
\externaldocument[spaces-over-fields-]{spaces-over-fields}
\externaldocument[spaces-topologies-]{spaces-topologies}
\externaldocument[spaces-descent-]{spaces-descent}
\externaldocument[spaces-perfect-]{spaces-perfect}
\externaldocument[spaces-more-morphisms-]{spaces-more-morphisms}
\externaldocument[spaces-flat-]{spaces-flat}
\externaldocument[spaces-groupoids-]{spaces-groupoids}
\externaldocument[spaces-more-groupoids-]{spaces-more-groupoids}
\externaldocument[bootstrap-]{bootstrap}
\externaldocument[spaces-pushouts-]{spaces-pushouts}
\externaldocument[spaces-chow-]{spaces-chow}
\externaldocument[groupoids-quotients-]{groupoids-quotients}
\externaldocument[spaces-more-cohomology-]{spaces-more-cohomology}
\externaldocument[spaces-simplicial-]{spaces-simplicial}
\externaldocument[spaces-duality-]{spaces-duality}
\externaldocument[formal-spaces-]{formal-spaces}
\externaldocument[restricted-]{restricted}
\externaldocument[spaces-resolve-]{spaces-resolve}
\externaldocument[formal-defos-]{formal-defos}
\externaldocument[defos-]{defos}
\externaldocument[cotangent-]{cotangent}
\externaldocument[examples-defos-]{examples-defos}
\externaldocument[algebraic-]{algebraic}
\externaldocument[examples-stacks-]{examples-stacks}
\externaldocument[stacks-sheaves-]{stacks-sheaves}
\externaldocument[criteria-]{criteria}
\externaldocument[artin-]{artin}
\externaldocument[quot-]{quot}
\externaldocument[stacks-properties-]{stacks-properties}
\externaldocument[stacks-morphisms-]{stacks-morphisms}
\externaldocument[stacks-limits-]{stacks-limits}
\externaldocument[stacks-cohomology-]{stacks-cohomology}
\externaldocument[stacks-perfect-]{stacks-perfect}
\externaldocument[stacks-introduction-]{stacks-introduction}
\externaldocument[stacks-more-morphisms-]{stacks-more-morphisms}
\externaldocument[stacks-geometry-]{stacks-geometry}
\externaldocument[moduli-]{moduli}
\externaldocument[moduli-curves-]{moduli-curves}
\externaldocument[examples-]{examples}
\externaldocument[exercises-]{exercises}
\externaldocument[guide-]{guide}
\externaldocument[desirables-]{desirables}
\externaldocument[coding-]{coding}
\externaldocument[obsolete-]{obsolete}
\externaldocument[fdl-]{fdl}
\externaldocument[index-]{index}

% Theorem environments.
%
\theoremstyle{plain}
\newtheorem{theorem}[subsection]{Theorem}
\newtheorem{proposition}[subsection]{Proposition}
\newtheorem{lemma}[subsection]{Lemma}

\theoremstyle{definition}
\newtheorem{definition}[subsection]{Definition}
\newtheorem{example}[subsection]{Example}
\newtheorem{exercise}[subsection]{Exercise}
\newtheorem{situation}[subsection]{Situation}

\theoremstyle{remark}
\newtheorem{remark}[subsection]{Remark}
\newtheorem{remarks}[subsection]{Remarks}

\numberwithin{equation}{subsection}

% Macros
%
\def\lim{\mathop{\mathrm{lim}}\nolimits}
\def\colim{\mathop{\mathrm{colim}}\nolimits}
\def\Spec{\mathop{\mathrm{Spec}}}
\def\Hom{\mathop{\mathrm{Hom}}\nolimits}
\def\Ext{\mathop{\mathrm{Ext}}\nolimits}
\def\SheafHom{\mathop{\mathcal{H}\!\mathit{om}}\nolimits}
\def\SheafExt{\mathop{\mathcal{E}\!\mathit{xt}}\nolimits}
\def\Sch{\mathit{Sch}}
\def\Mor{\mathop{\mathrm{Mor}}\nolimits}
\def\Ob{\mathop{\mathrm{Ob}}\nolimits}
\def\Sh{\mathop{\mathit{Sh}}\nolimits}
\def\NL{\mathop{N\!L}\nolimits}
\def\CH{\mathop{\mathrm{CH}}\nolimits}
\def\proetale{{pro\text{-}\acute{e}tale}}
\def\etale{{\acute{e}tale}}
\def\QCoh{\mathit{QCoh}}
\def\Ker{\mathop{\mathrm{Ker}}}
\def\Im{\mathop{\mathrm{Im}}}
\def\Coker{\mathop{\mathrm{Coker}}}
\def\Coim{\mathop{\mathrm{Coim}}}

% Boxtimes
%
\DeclareMathSymbol{\boxtimes}{\mathbin}{AMSa}{"02}

%
% Macros for moduli stacks/spaces
%
\def\QCohstack{\mathcal{QC}\!\mathit{oh}}
\def\Cohstack{\mathcal{C}\!\mathit{oh}}
\def\Spacesstack{\mathcal{S}\!\mathit{paces}}
\def\Quotfunctor{\mathrm{Quot}}
\def\Hilbfunctor{\mathrm{Hilb}}
\def\Curvesstack{\mathcal{C}\!\mathit{urves}}
\def\Polarizedstack{\mathcal{P}\!\mathit{olarized}}
\def\Complexesstack{\mathcal{C}\!\mathit{omplexes}}
% \Pic is the operator that assigns to X its picard group, usage \Pic(X)
% \Picardstack_{X/B} denotes the Picard stack of X over B
% \Picardfunctor_{X/B} denotes the Picard functor of X over B
\def\Pic{\mathop{\mathrm{Pic}}\nolimits}
\def\Picardstack{\mathcal{P}\!\mathit{ic}}
\def\Picardfunctor{\mathrm{Pic}}
\def\Deformationcategory{\mathcal{D}\!\mathit{ef}}
\usepackage{cleveref}
\crefname{originallemma}{original lemma}{original lemmas}
\Crefname{originallemma}{Original lemma}{Original lemmas}
\makeatletter
\newlabel{homology-lemma-additive-cat-biproduct-kernel}{{09QG}{source}{}{sourcepointer-1}{}}
\newlabel{homology-lemma-additive-cat-biproduct-kernel@cref}{{[originallemma][0][]09QG}{[1][0][]source}{}{}{}}
\makeatother
\title{The dg analogue of the octahedral axiom TR4 for two composable admissible monomorphisms}
\author{The Stacks Project Authors}
\date{stacks-project:dga.tex:local-snapshot-20261012; GNU FDL 1.2 or later}
\begin{document}
\maketitle
\noindent\textit{Editorial note.} The counted statement is the author's own lemma environment labelled lemma-dgc-analogue-tr4 (source lines 4576-4583); the proof environment begins on the second line after it and is its own proof, closing at source line 4688.\par\medskip
\noindent\textit{Editorial note.} The statement and the proof invoke the author's Situation ABC and the triangle lemma of the same chapter; both are retained as context spans and the generator expands the context transitively (budget 1200 lines for this card, since the triangle lemma drags in the whole splitting apparatus), while reference targets in other chapters that carry canonical project tags (here the Homology biproduct-kernel lemma) are auto-declared as non-load-bearing source pointers.\par\medskip
\noindent\textit{Editorial note.} The retained preamble is the chapter's own minimal preamble (sources/stacks-dga/preamble.tex, taken with its real line count), a verbatim excerpt of the native Xy-pic setup, the native theorem declarations with the mathematical macros, and the project's own no-op reference environments; no macro was invented or redefined.\par\medskip
\noindent\textit{Editorial scope.} One original result of the Stacks Project's Differential Graded Algebra chapter, counted once: the octahedral axiom TR4 holds for two composable admissible monomorphisms in the author's Situation ABC.\par\medskip
\noindent\textit{Allowed standard prerequisites (editorial).}\begin{itemize}
\item Situation ABC. The author's standing situation: a differential graded category with admissible monomorphisms, admissible short exact sequences and the associated graded additive category.. It is the setting of the statement and of axiom (A) used at source line 4655 (source line 4578).
\item Distinguished triangles from admissible short exact sequences. Every admissible short exact sequence extends to a distinguished triangle, with the connecting map defined as pi d(s) for chosen splittings, and the resulting triangles satisfy the triangle identities (the author's lemma-get-triangle, retained as context).. It produces the three triangles and supplies every computational identity used in the proof (source lines 4587-4615, 4636-4644, 4674-4675).
\item Biproduct and kernel in an additive category. In an additive category, the kernel of the projection of a biproduct onto its first factor can be computed as the complementary summand (Homology, Lemma homology-lemma-additive-cat-biproduct-kernel, declared as a tagged source pointer).. It gives q\_2 = q\_1 direct sum q\_3 in the underlying graded category (source lines 4657-4666).
\item The octahedral axiom TR4. The statement of TR4 for two composable morphisms in a triangulated category.. It is what the constructed triangles are asserted to satisfy (source line 4582).
\end{itemize}\medskip
\ifdefined\hypertarget\hypertarget{sourcepointer-1}{}\fi
\noindent\textit{Source reference.} Original lemma 09QG: sources/stacks-chow/tags.txt, source lines 12284--12284. Cross-chapter pointer whose target carries the canonical tag 09QG in the project tag list; declared so the reference resolves inside the assembled package, while every step the counted argument uses remains inside the retained spans.\par\medskip
\begin{situation}
\label{situation-ABC}
Here $R$ is a ring and $\mathcal{A}$ is a differential graded category
over $R$ having axioms (A), (B), and
\begin{enumerate}
\item[(C)] given an arrow $f : x \to y$ of degree $0$ with
$\text{d}(f) = 0$ there exists an admissible short exact sequence
$y \to c(f) \to x[1]$ in $\text{Comp}(\mathcal{A})$ such that the map
$x[1] \to y[1]$ of Lemma \ref{lemma-get-triangle} is equal to $f[1]$.
\end{enumerate}
\end{situation}

\begin{lemma}
\label{lemma-get-triangle}
Let $\mathcal{A}$ be a differential graded category satisfying
axioms (A) and (B). Given an admissible short exact sequence
$x \to y \to z$ we obtain (see proof) a triangle
$$
x \to y \to z \to x[1]
$$
in $\text{Comp}(\mathcal{A})$ with the property that any two compositions
in $z[-1] \to x \to y \to z \to x[1]$ are zero in $K(\mathcal{A})$.
\end{lemma}

\begin{proof}
Choose a diagram
$$
\xymatrix{
x \ar[rr]_1 \ar[rd]_a & & x \\
& y \ar[ru]_\pi \ar[rd]^b & \\
z \ar[rr]^1 \ar[ru]^s & & z
}
$$
giving the isomorphism of graded objects $y \cong x \oplus z$ as in the
definition of an admissible short exact sequence. Here are some equations
that hold in this situation
\begin{enumerate}
\item $1 = \pi a$ and hence $\text{d}(\pi) a = 0$,
\item $1 = b s$ and hence $b \text{d}(s) = 0$,
\item $1 = a \pi + s b$ and hence $a \text{d}(\pi) + \text{d}(s) b = 0$,
\item $\pi s = 0$ and hence $\text{d}(\pi)s + \pi \text{d}(s) = 0$,
\item $\text{d}(s) = a \pi \text{d}(s)$ because
$\text{d}(s) = (a \pi + s b)\text{d}(s)$ and $b\text{d}(s) = 0$,
\item $\text{d}(\pi) = \text{d}(\pi) s b$ because
$\text{d}(\pi) = \text{d}(\pi)(a \pi + s b)$ and $\text{d}(\pi)a = 0$,
\item $\text{d}(\pi \text{d}(s)) = 0$ because if we postcompose it
with the monomorphism $a$ we get
$\text{d}(a\pi \text{d}(s)) = \text{d}(\text{d}(s)) = 0$, and
\item $\text{d}(\text{d}(\pi)s) = 0$ as by (4) it is the negative
of $\text{d}(\pi\text{d}(s))$ which is $0$ by (7).
\end{enumerate}
We've used repeatedly that $\text{d}(a) = 0$, $\text{d}(b) = 0$,
and that $\text{d}(1) = 0$. By (7) we see that
$$
\delta = \pi \text{d}(s) = - \text{d}(\pi) s : z \to x[1]
$$
is a morphism in $\text{Comp}(\mathcal{A})$. By (5) we see that
the composition $a \delta = a \pi \text{d}(s) = \text{d}(s)$
is homotopic to zero. By (6) we see that the composition
$\delta b = - \text{d}(\pi)sb = \text{d}(-\pi)$ is homotopic to zero.
\end{proof}

\begin{lemma}
\label{lemma-dgc-analogue-tr4}
In Situation \ref{situation-ABC} given admissible monomorphisms
$x \xrightarrow{\alpha} y$, $y \xrightarrow{\beta} z$ in $\mathcal{A}$,
there exist distinguished triangles
$(x,y,q_1,\alpha,p_1,\delta_1)$, $(x,z,q_2,\beta\alpha,p_2,\delta_2)$
and $(y,z,q_3,\beta,p_3,\delta_3)$ for which TR4 holds.
\end{lemma}

\begin{proof}
Given admissible monomorphisms $x\xrightarrow{\alpha} y$ and
$y\xrightarrow{\beta}z$, we can find distinguished triangles,
via their extensions to admissible short exact sequences,
$$
\xymatrix{
x\ar@<0.5ex>[r]^{\alpha} &
y\ar@<0.5ex>[l]^{\pi_1} \ar@<0.5ex>[r]^{p_1} &
q_1 \ar[r]^{\delta_1} \ar@<0.5ex>[l]^{s_1} &
x[1]
}
$$
$$
\xymatrix{
x\ar@<0.5ex>[r]^{\beta\alpha} &
z\ar@<0.5ex>[l]^{\pi_1\pi_3} \ar@<0.5ex>[r]^{p_2} &
q_2 \ar[r]^{\delta_2} \ar@<0.5ex>[l]^{s_2} &
x[1]
}
$$
$$
\xymatrix{
y\ar@<0.5ex>[r]^{\beta} &
z\ar@<0.5ex>[l]^{\pi_3} \ar@<0.5ex>[r]^{p_3} &
q_3 \ar[r]^{\delta_3} \ar@<0.5ex>[l]^{s_3} &
x[1]
}
$$
In these diagrams, the maps $\delta_i$ are defined as
$\delta_i = \pi_i d(s_i)$ analogous to the maps defined in
Lemma \ref{lemma-get-triangle}.
They fit in the following solid commutative diagram
$$
\xymatrix@C=5pc@R=3pc{
x\ar@<0.5ex>[r]^{\alpha} \ar@<0.5ex>[dr]^{\beta\alpha} &
y\ar@<0.5ex>[d]^{\beta} \ar@<0.5ex>[l]^{\pi_1} \ar@<0.5ex>[r]^{p_1} &
q_1 \ar[r]^{\delta_1} \ar@<0.5ex>[l]^{s_1} \ar@{.>}[dd]^{p_2\beta s_1} &
x[1] \\
 &
z \ar@<0.5ex>[u]^{\pi_3}\ar@<0.5ex>[d]^{p_3}
\ar@<0.5ex>[dr]^{p_2} \ar@<0.5ex>[ul]^{\pi_1\pi_3} & & \\
 &
q_3\ar@<0.5ex>[u]^{s_3} \ar[d]^{\delta_3} &
q_2 \ar@{.>}[l]^{p_3s_2} \ar@<0.5ex>[ul]^{s_2} \ar[dr]^{\delta_2} \\
 &
y[1] & & x[1]}
$$
where we have defined the dashed arrows as indicated.
Clearly, their composition $p_3s_2p_2\beta s_1 = 0$
since $s_2p_2 = 0$. We claim that they both are morphisms of
$\text{Comp}(\mathcal{A})$. We can check this using equations in
Lemma \ref{lemma-get-triangle}:
$$
d(p_2\beta s_1) = p_2\beta d(s_1) = p_2\beta\alpha\pi_1 d(s_1) = 0
$$
since $p_2\beta\alpha = 0$, and
$$
d(p_3s_2) = p_3d(s_2) = p_3\beta\alpha\pi_1\pi_3 d(s_2) = 0
$$
since $p_3\beta = 0$. To check that $q_1\to q_2\to q_3$
is an admissible short exact sequence, it remains to show
that in the underlying graded category, $q_2 = q_1\oplus q_3$
with the above two morphisms as coprojection and projection.
To do this, observe that in the underlying graded category
$\mathcal{C}$, there hold
$$
y = x\oplus q_1,\quad
z = y\oplus q_3 = x\oplus q_1\oplus q_3
$$
where $\pi_1\pi_3$ gives the projection morphism onto the first
factor: $x\oplus q_1\oplus q_3\to z$. By axiom (A) on
$\mathcal{A}$, $\mathcal{C}$ is an additive category, hence
we may apply
Homology, Lemma \ref{homology-lemma-additive-cat-biproduct-kernel}
and conclude that
$$
\Ker(\pi_1\pi_3) = q_1\oplus q_3
$$
in $\mathcal{C}$. Another application of
Homology, Lemma \ref{homology-lemma-additive-cat-biproduct-kernel}
to $z = x\oplus q_2$ gives $\Ker(\pi_1\pi_3) = q_2$.
Hence $q_2\cong q_1\oplus q_3$ in $\mathcal{C}$.
It is clear that the dashed morphisms defined above give
coprojection and projection.

\medskip\noindent
Finally, we have to check that the morphism
$\delta : q_3 \to q_1[1]$ induced by the admissible
short exact sequence $q_1\to q_2\to q_3$ agrees with
$p_1\delta_3$. By the construction in
Lemma \ref{lemma-get-triangle}, the morphism $\delta$ is given by
\begin{align*}
p_1\pi_3s_2d(p_2s_3)
= &
p_1\pi_3s_2p_2d(s_3) \\
= &
p_1\pi_3(1-\beta\alpha\pi_1\pi_3)d(s_3) \\
= &
p_1\pi_3d(s_3)\quad (\text{since }\pi_3\beta = 0) \\
= &
p_1\delta_3
\end{align*}
as desired. The proof is complete.
\end{proof}

\end{document}
